\documentclass[11pt,a4paper]{article}

% Packages
\usepackage[utf8]{inputenc}
\usepackage[T1]{fontenc}
\usepackage{amsmath,amssymb,amsfonts}
\usepackage[margin=1in]{geometry}
\usepackage{enumitem}
\usepackage{hyperref}
\usepackage{fancyhdr}
\usepackage{titlesec}
\usepackage{tocloft}

% Page style
\pagestyle{fancy}
\fancyhf{}
\fancyhead[L]{\small Lecture Notes: Functions}
\fancyhead[R]{\small \thepage}
\renewcommand{\headrulewidth}{0.4pt}

% Section formatting
\titleformat{\section}{\Large\bfseries}{Section \thesection}{1em}{}
\titleformat{\subsection}{\large\bfseries}{\thesubsection}{1em}{}

% Adjust table of contents appearance
\renewcommand{\cftsecfont}{\bfseries}
\renewcommand{\cftsubsecfont}{\normalfont}

\begin{document}

% Title page
\begin{titlepage}
    \centering
    \vspace*{4cm}
    
    {\Huge\bfseries Lecture Notes:\\ Functions\par}
    \vspace{1cm}
    {\Large Discrete Mathematics\par}
    \vspace{2cm}
    {\large \textbf{Author:} Bakdaulet Sotsial\par}
    \vspace{1cm}
    
    \vfill
    {\small Astana IT University \par}
    \vspace*{2cm}
\end{titlepage}

% Table of contents
\tableofcontents
\newpage

% Section 1: Functions
\section{Functions}

\textbf{Definition (Function).} Let $A$ and $B$ be non-empty sets. A \textit{function} (or \textit{mapping}) $f$ from $A$ to $B$, denoted $f: A \to B$, is a rule that assigns to each element $x \in A$ exactly one element $y \in B$. The element $y$ is called the \textit{image} of $x$ under $f$ and is written $f(x)$.

\textbf{Remark.} A function is a special type of binary relation: it is a relation $f \subseteq A \times B$ such that for every $a \in A$, there exists exactly one $b \in B$ with $(a, b) \in f$. This is sometimes called a "well-defined" assignment.

\subsection{Domain, Codomain, and Range}

\begin{itemize}[label=\(\triangleright\)]
    \item \textbf{Domain:} The set $A$ is called the \textit{domain} of $f$, denoted $\operatorname{dom}(f)$. It is the set of all possible inputs.
    \item \textbf{Codomain:} The set $B$ is called the \textit{codomain} of $f$. It is the set within which all outputs lie.
    \item \textbf{Range (Image):} The \textit{range} or \textit{image} of $f$, denoted $\operatorname{ran}(f)$ or $f(A)$, is the set of all actual outputs:
    \begin{equation*}
        \operatorname{ran}(f) = \{f(x) \mid x \in A\} \subseteq B.
    \end{equation*}
\end{itemize}

\textbf{Remark.} Note carefully: the codomain is part of the definition of the function, while the range is a derived set that is always a subset of the codomain.

\begin{example}
Let $f: \mathbb{R} \to \mathbb{R}$ be defined by $f(x) = x^2$.
\begin{itemize}[label=--]
    \item Domain: $\mathbb{R}$.
    \item Codomain: $\mathbb{R}$.
    \item Range: $[0, \infty) = \{y \in \mathbb{R} \mid y \geq 0\}$.
\end{itemize}
Notice that the range is a proper subset of the codomain.
\end{example}

\begin{example}
Consider $g: \{1, 2, 3\} \to \{a, b, c, d\}$ defined by
\begin{equation*}
    g(1) = a, \quad g(2) = c, \quad g(3) = a.
\end{equation*}
Then $\operatorname{dom}(g) = \{1, 2, 3\}$, codomain is $\{a, b, c, d\}$, and $\operatorname{ran}(g) = \{a, c\}$.
\end{example}

\subsection{Function Notation and Equality}

We write $f: A \to B$ to denote a function, and for $x \in A$, the output is denoted $f(x)$.

\textbf{Definition (Equality of Functions).} Two functions $f: A \to B$ and $g: C \to D$ are \textit{equal}, written $f = g$, if and only if:
\begin{enumerate}[label=(\roman*)]
    \item $A = C$ (same domain),
    \item $B = D$ (same codomain),
    \item $f(x) = g(x)$ for all $x \in A$ (same rule).
\end{enumerate}

\textbf{Definition (Image of a Set).} For a subset $X \subseteq A$, the \textit{image} of $X$ under $f$ is
\begin{equation*}
    f(X) = \{f(x) \mid x \in X\}.
\end{equation*}
For a subset $Y \subseteq B$, the \textit{preimage} (or \textit{inverse image}) of $Y$ is
\begin{equation*}
    f^{-1}(Y) = \{x \in A \mid f(x) \in Y\}.
\end{equation*}

% Section 2: Types of Functions
\section{Types of Functions}

Functions can be classified based on how they map elements from the domain to the codomain. The three most important classes are injective, surjective, and bijective functions.

\subsection{One-to-One (Injective) Functions}

\textbf{Definition (Injective Function).} A function $f: A \to B$ is \textit{injective} (or \textit{one-to-one}) if distinct elements of the domain map to distinct elements of the codomain. Formally:
\begin{equation*}
    \forall x_1, x_2 \in A,\; \bigl(f(x_1) = f(x_2) \implies x_1 = x_2\bigr).
\end{equation*}
Equivalently (contrapositive): $x_1 \neq x_2 \implies f(x_1) \neq f(x_2)$.

\textbf{Intuition:} In an injective function, no two different inputs produce the same output. Each element of the range is hit by at most one element from the domain.

\begin{example}[Injective]
The function $f: \mathbb{R} \to \mathbb{R}$ given by $f(x) = 3x + 2$ is injective. Indeed, if $3x_1 + 2 = 3x_2 + 2$, then $3x_1 = 3x_2$, so $x_1 = x_2$.
\end{example}

\begin{example}[Not Injective]
The function $g: \mathbb{R} \to \mathbb{R}$ given by $g(x) = x^2$ is \textbf{not} injective, because $g(2) = 4$ and $g(-2) = 4$, yet $2 \neq -2$.
\end{example}

\subsection{Onto (Surjective) Functions}

\textbf{Definition (Surjective Function).} A function $f: A \to B$ is \textit{surjective} (or \textit{onto}) if every element of the codomain $B$ is the image of at least one element of the domain $A$. Formally:
\begin{equation*}
    \forall y \in B,\; \exists x \in A \text{ such that } f(x) = y.
\end{equation*}
Equivalently, the range equals the codomain: $\operatorname{ran}(f) = B$.

\textbf{Intuition:} A surjective function "covers" the entire codomain: every element in $B$ is "hit" by the function.

\begin{example}[Surjective]
The function $f: \mathbb{R} \to \mathbb{R}$ given by $f(x) = x^3$ is surjective. For any $y \in \mathbb{R}$, we can take $x = \sqrt[3]{y}$ (a real number) and $f(x) = y$.
\end{example}

\begin{example}[Not Surjective]
The function $g: \mathbb{R} \to \mathbb{R}$ given by $g(x) = x^2$ is \textbf{not} surjective, because there is no real $x$ such that $g(x) = -1$ (the range is only $[0, \infty)$).
\end{example}

\subsection{Bijective Functions}

\textbf{Definition (Bijective Function).} A function $f: A \to B$ is \textit{bijective} (or a \textit{one-to-one correspondence}) if it is both injective and surjective. That is, every element of $B$ is the image of exactly one element of $A$.

\textbf{Intuition:} A bijection pairs up elements of $A$ and $B$ uniquely in both directions. This is the key concept for comparing the "sizes" of sets.

\begin{example}[Bijective]
The function $f: \mathbb{R} \to \mathbb{R}$ given by $f(x) = 2x + 1$ is bijective.
\begin{itemize}[label=--]
    \item Injective: $2x_1 + 1 = 2x_2 + 1 \implies x_1 = x_2$.
    \item Surjective: For any $y \in \mathbb{R}$, $x = \frac{y-1}{2}$ satisfies $f(x) = y$.
\end{itemize}
\end{example}

\begin{example}[Bijective on Restricted Domain]
The function $h: [0, \infty) \to [0, \infty)$ defined by $h(x) = x^2$ is bijective. It is injective because on non-negative numbers, $x_1^2 = x_2^2$ implies $x_1 = x_2$; it is surjective because every non-negative $y$ has a non-negative square root $x = \sqrt{y}$.
\end{example}

% Section 3: Inverse Functions
\section{Inverse Functions}

\textbf{Definition (Inverse Function).} Let $f: A \to B$ be a bijective function. The \textit{inverse function} of $f$, denoted $f^{-1}: B \to A$, is the function that assigns to each $y \in B$ the unique $x \in A$ such that $f(x) = y$. That is,
\begin{equation*}
    f^{-1}(y) = x \quad \text{if and only if} \quad f(x) = y.
\end{equation*}

\textbf{Theorem (Existence of Inverse).} A function $f: A \to B$ has an inverse function $f^{-1}: B \to A$ if and only if $f$ is bijective.

\textbf{Proof Sketch.} If $f$ is not injective, then $f(x_1) = f(x_2) = y$ for some $x_1 \neq x_2$, making it impossible to define $f^{-1}(y)$ uniquely. If $f$ is not surjective, then some $y \in B$ has no preimage, so $f^{-1}(y)$ is undefined. Conversely, if $f$ is bijective, each $y$ has exactly one $x$ with $f(x) = y$, defining a proper function.

\subsection{Finding Inverse Functions}

To find the inverse of a bijective function $f$, solve the equation $y = f(x)$ for $x$ in terms of $y$, then interchange the roles of $x$ and $y$ (or simply express $f^{-1}(x)$).

\begin{example}
Let $f: \mathbb{R} \to \mathbb{R}$ be defined by $f(x) = 3x - 7$. Find $f^{-1}$.

\textbf{Solution.} Set $y = 3x - 7$. Solve for $x$:
\begin{align*}
    y &= 3x - 7 \\
    y + 7 &= 3x \\
    x &= \frac{y + 7}{3}.
\end{align*}
Thus, $f^{-1}(y) = \frac{y+7}{3}$, or equivalently, $f^{-1}(x) = \frac{x+7}{3}$.
\end{example}

\begin{example}
Let $g: [0, \infty) \to [0, \infty)$ be $g(x) = x^2 + 1$. Find $g^{-1}$.

\textbf{Solution.} Set $y = x^2 + 1$. Solve for $x$ with $x \geq 0$:
\begin{align*}
    y - 1 &= x^2 \\
    x &= \sqrt{y - 1} \quad (\text{taking the non-negative root since } x \geq 0).
\end{align*}
Thus, $g^{-1}(x) = \sqrt{x - 1}$, with domain $[1, \infty)$ (which matches $\operatorname{ran}(g)$).
\end{example}

\subsection{Verification by Composition}

\textbf{Theorem (Cancellation Laws).} If $f: A \to B$ is bijective and $f^{-1}: B \to A$ is its inverse, then:
\begin{align*}
    f^{-1}(f(x)) &= x \quad \text{for all } x \in A, \\
    f(f^{-1}(y)) &= y \quad \text{for all } y \in B.
\end{align*}

This provides a practical way to verify an inverse: compose the functions in both orders and check that the identity results.

\begin{example}
For $f(x) = 3x - 7$ and $f^{-1}(x) = \frac{x+7}{3}$:
\begin{align*}
    f^{-1}(f(x)) &= f^{-1}(3x - 7) = \frac{(3x - 7) + 7}{3} = \frac{3x}{3} = x, \\
    f(f^{-1}(x)) &= f\!\left(\frac{x+7}{3}\right) = 3\left(\frac{x+7}{3}\right) - 7 = x + 7 - 7 = x.
\end{align*}
\end{example}

% Section 4: Composition of Functions
\section{Composition of Functions}

\textbf{Definition (Composition).} Let $f: A \to B$ and $g: B \to C$ be functions. The \textit{composition} of $g$ with $f$, denoted $g \circ f$ (read "g after f" or "g of f"), is the function $g \circ f: A \to C$ defined by
\begin{equation*}
    (g \circ f)(x) = g(f(x)) \quad \text{for all } x \in A.
\end{equation*}

\textbf{Remark (Domain Considerations).} For $g \circ f$ to be defined, the codomain of $f$ must equal (or be a subset of) the domain of $g$. The diagram is:
\[
\begin{array}{ccccc}
    A & \xrightarrow{f} & B & \xrightarrow{g} & C \\
    \multicolumn{5}{c}{A \xrightarrow{g \circ f} C}
\end{array}
\]

\subsection{Properties of Composition}

\textbf{Theorem (Associativity).} Composition of functions is associative. If $f: A \to B$, $g: B \to C$, and $h: C \to D$, then
\begin{equation*}
    (h \circ g) \circ f = h \circ (g \circ f).
\end{equation*}

\textbf{Proof.} For any $x \in A$:
\begin{align*}
    ((h \circ g) \circ f)(x) &= (h \circ g)(f(x)) = h(g(f(x))) \\
    (h \circ (g \circ f))(x) &= h((g \circ f)(x)) = h(g(f(x))).
\end{align*}
Both are equal, so associativity holds.

\textbf{Remark (Non-Commutativity).} Composition of functions is \textbf{not} commutative in general. Even when both orders are defined, $f \circ g$ and $g \circ f$ are usually different.

\begin{example}
Let $f, g: \mathbb{R} \to \mathbb{R}$ be given by $f(x) = x^2$ and $g(x) = x + 1$. Then:
\begin{align*}
    (f \circ g)(x) &= f(g(x)) = f(x+1) = (x+1)^2 = x^2 + 2x + 1, \\
    (g \circ f)(x) &= g(f(x)) = g(x^2) = x^2 + 1.
\end{align*}
Clearly, $(f \circ g)(x) \neq (g \circ f)(x)$ (for example, evaluate at $x = 0$).
\end{example}

\subsection{Composition with Inverse}

\textbf{Theorem.} If $f: A \to B$ is bijective, then
\begin{align*}
    f^{-1} \circ f &= \operatorname{id}_A \quad \text{(the identity function on } A\text{)}, \\
    f \circ f^{-1} &= \operatorname{id}_B \quad \text{(the identity function on } B\text{)}.
\end{align*}

This property is often used to prove that a function is bijective by constructing its inverse.

\subsection{Composition and Function Types}

\textbf{Theorem.} Let $f: A \to B$ and $g: B \to C$.
\begin{enumerate}[label=(\roman*)]
    \item If $f$ and $g$ are both injective, then $g \circ f$ is injective.
    \item If $f$ and $g$ are both surjective, then $g \circ f$ is surjective.
    \item If $f$ and $g$ are both bijective, then $g \circ f$ is bijective.
\end{enumerate}

% Section 5: Cardinality of Sets
\section{Cardinality of Sets}

\textbf{Definition (Cardinality).} The \textit{cardinality} of a set $A$, denoted $|A|$, is a measure of the "number of elements" in $A$. For finite sets, it is simply the count of distinct elements.

\textbf{Definition (Finite and Infinite Sets).}
\begin{itemize}
    \item A set is \textit{finite} if it has $n$ elements for some $n \in \mathbb{N} \cup \{0\}$; then $|A| = n$.
    \item A set is \textit{infinite} if it is not finite.
\end{itemize}

\subsection{Comparing Cardinalities via Functions}

Functions, especially bijections, provide a rigorous way to compare the sizes of sets without counting elements explicitly.

\textbf{Definition (Equinumerosity).} Two sets $A$ and $B$ have the same \textit{cardinality} (are \textit{equinumerous} or \textit{equipotent}), denoted $|A| = |B|$, if there exists a bijection $f: A \to B$.

\textbf{Theorem (Cantor—Bernstein—Schröder).} If there exists an injection $f: A \to B$ and an injection $g: B \to A$, then there exists a bijection $h: A \to B$, so $|A| = |B|$. (We state this without proof here.)

\begin{example}
Let $A = \{1, 2, 3, 4\}$ and $B = \{a, b, c, d\}$. The function $f: A \to B$ defined by $f(1)=a, f(2)=b, f(3)=c, f(4)=d$ is a bijection, so $|A| = |B| = 4$.
\end{example}

\begin{example}[Even and Natural Numbers]
Let $\mathbb{N} = \{1, 2, 3, \dots\}$ and $E = \{2, 4, 6, \dots\}$. The function $f: \mathbb{N} \to E$ defined by $f(n) = 2n$ is a bijection. Therefore, $|\mathbb{N}| = |E|$, even though $E \subset \mathbb{N}$. This demonstrates that an infinite set can have the same cardinality as a proper subset of itself — a defining characteristic of infinite sets (Dedekind's definition).
\end{example}

\subsection{Countable and Uncountable Sets}

\textbf{Definition (Countable Set).} A set $A$ is \textit{countable} if it is finite or if there exists a bijection $f: \mathbb{N} \to A$ (equivalently, $|A| = |\mathbb{N}|$, denoted $\aleph_0$, "aleph-null"). In this case, the elements can be listed as $a_1, a_2, a_3, \dots$

\textbf{Examples of Countable Sets:}
\begin{itemize}[label=--]
    \item $\mathbb{N}$, $\mathbb{Z}$, $\mathbb{Q}$ (rational numbers) are countably infinite.
    \item Any finite set is countable by definition.
\end{itemize}

\textbf{Theorem (Cantor's Diagonal Argument).} The set of real numbers $\mathbb{R}$ is \textit{uncountable}. That is, there is no bijection from $\mathbb{N}$ to $\mathbb{R}$. The cardinality of $\mathbb{R}$ is strictly greater than $\aleph_0$, often denoted $\mathfrak{c}$ (the continuum).

\begin{example}[Countability of $\mathbb{Z}$]
We can show $|\mathbb{Z}| = |\mathbb{N}|$ by constructing an explicit bijection $f: \mathbb{N} \to \mathbb{Z}$:
\begin{equation*}
    f(n) = \begin{cases}
        \frac{n}{2} & \text{if } n \text{ is even}, \\[4pt]
        -\frac{n-1}{2} & \text{if } n \text{ is odd}.
    \end{cases}
\end{equation*}
The listing order is: $0, 1, -1, 2, -2, 3, -3, \dots$ (if we consider $\mathbb{Z}$ including 0, with $\mathbb{N}$ shifted appropriately). This demonstrates that $\mathbb{Z}$ is countably infinite.
\end{example}

\begin{example}[Uncountability of the Interval $(0, 1)$]
The open interval $(0, 1)$ is uncountable. Any attempted list $r_1, r_2, r_3, \dots$ of real numbers in $(0, 1)$, expressed in decimal expansion, can be shown to miss some number by constructing a new real number that differs from $r_k$ in the $k$-th decimal place (Cantor's diagonal argument). Hence, $(0,1)$ is not countable.
\end{example}

\textbf{Remark.} The concept of cardinality is a profound extension of counting to infinite sets. The study of infinite cardinalities leads to the hierarchy:
\begin{equation*}
    |\mathbb{N}| = \aleph_0 < |\mathbb{R}| = \mathfrak{c} < |\mathcal{P}(\mathbb{R})| < \cdots
\end{equation*}
A fundamental question (the Continuum Hypothesis) asks whether there exists a set with cardinality strictly between $\aleph_0$ and $\mathfrak{c}$.

\end{document}